\chapter{Systeemarchitectuur en werkstroomdiagrammen}\label{app:architectuur}

Figuur~\ref{fig:takaro-mcp-loop} toont de volledige architectuur van de eval-runner en orkestratieworkflow. Het diagram illustreert hoe de eval-runner model-promptcombinaties itereert, per combinatie een Claude Code-instantie spawnt, en elke sessie via een \texttt{/eval-hint} koppelt aan de MCP-server en Langfuse voor gestructureerde tracing. De agentische lus toont de systeempromptbranching tussen standaard- en orchestratormodus, de toolcall-cyclus en de opbouw van evaluatiemetrieken.

\begin{figure}[p]
    \includegraphics[width=\textwidth]{TakaroMCPArchitecture.png}
    \centering
    \caption[Agentic loop en MCP]{Volledige architectuur van de eval-runner en orkestratieworkflow. De eval-runner itereert over model-promptcombinaties, spawnt per combinatie een Claude Code-instantie en koppelt elke sessie via \texttt{/eval-hint} aan de MCP-server voor gestructureerde Langfuse-tracing.}\label{fig:takaro-mcp-loop}
\end{figure}

\begin{figure}[p]
    \includegraphics[scale=0.70]{scaffoldModule.png}
    \centering
    \caption[Take Action fase]{Take Action fase. Hier wordt besilst op welke manier de code voor de module gegenereert zal worden.}\label{fig:takaro-mcp-delegate}
\end{figure}

\begin{figure}[p]
    \includegraphics[scale=0.70]{pushModule.png}
    \centering
    \caption[Push module workflow]{Vervolg van de take action fase. Hier word de workflow doornomen na het pushen een instaleren van een module op de takaro server.}\label{fig:takaro-mcp-push}
\end{figure}
